\chapter{SYSTEM ANALYSIS AND REQUIREMENTS} \label{c3}
\section{Existing System Analysis}
In existing classroom practices, teachers either call roll numbers manually or use basic attendance apps where students submit a simple check-in. These workflows do not guarantee that the correct student is physically present in class. They also require manual monitoring effort by faculty.

\section{Proposed System Analysis}
The proposed system introduces a constrained attendance workflow:
\begin{enumerate}
\item Teacher starts a session for selected date and period.
\item System generates a one-time code valid during active session.
\item Student enters code and validates session.
\item Student passes Bluetooth proximity and face verification.
\item Attendance is stored only if all validation checks succeed.
\end{enumerate}

\section{Functional Requirements}
\subsection{Teacher-side Requirements}
\begin{itemize}
\item Authenticate teacher and open dashboard.
\item Start session with class date and period.
\item Generate and display OTP with timer.
\item End active session manually or on expiry.
\item View live student list for ongoing session.
\item View slot-wise attendance and student reports.
\end{itemize}

\subsection{Student-side Requirements}
\begin{itemize}
\item Authenticate student and open dashboard.
\item Select class date and period.
\item Enter session OTP and validate active slot.
\item Perform Bluetooth check and face verification.
\item Submit attendance and view result status.
\end{itemize}

\subsection{Backend Requirements}
\begin{itemize}
\item Store user profiles, sessions, and attendance records.
\item Validate active sessions and enforce expiry.
\item Prevent duplicate attendance for the same student and slot.
\item Support teacher and student query operations.
\end{itemize}

\section{Non-Functional Requirements}
\begin{itemize}
\item \textbf{Usability:} The app must provide clear, guided steps with minimum confusion.
\item \textbf{Performance:} Session validation and submission should complete in practical classroom time.
\item \textbf{Reliability:} Retry and timeout handling should avoid app freezing behavior.
\item \textbf{Security:} Only authenticated requests should access attendance data.
\item \textbf{Maintainability:} Modular code and clean architecture for future upgrades.
\end{itemize}

\section{Feasibility Analysis}
\subsection{Technical Feasibility}
The project is technically feasible with Android SDK, Kotlin, BLE APIs, camera-based verification pipelines, and Supabase backend integration.

\subsection{Operational Feasibility}
The workflow is suitable for real classroom operation because it minimizes teacher actions and keeps student marking guided and controlled.

\subsection{Economic Feasibility}
The solution uses commonly available Android devices and free/low-cost development stack tools, making it economical for academic adoption.

\section{Requirement Summary Table}
\begin{table}[H]
\centering
\caption{Requirement Summary}
\label{c3:tab1}
\begin{tabular}{|L{3cm}|L{3.5cm}|L{5.5cm}|}
\hline
\textbf{Type} & \textbf{Module} & \textbf{Requirement}\\
\hline
Functional & Teacher & Start/stop session, OTP generation, live roster\\
\hline
Functional & Student & OTP validation, BLE and face verification, mark attendance\\
\hline
Functional & Backend & Persist profiles/sessions/attendance, duplicate prevention\\
\hline
Non-functional & UI/UX & Guided flow, status visibility, low confusion\\
\hline
Non-functional & Security & Authenticated access and verified attendance insertion\\
\hline
\end{tabular}
\end{table}

\section{Conclusion}
System analysis confirms that a multi-factor, session-controlled attendance system is both necessary and feasible for improving classroom attendance integrity.

